\documentclass{beamer}

% ===============================
% Theme and Packages
% ===============================
\usetheme{Madrid}
\usecolortheme{default}
\usefonttheme{professionalfonts}

\usepackage{amsmath, amssymb}
\usepackage{graphicx}
\usepackage{booktabs}
\usepackage{hyperref}
\usepackage{setspace}
\usepackage{ragged2e}

\begin{document}
	
	%------------------------------------------------
	% Slide 1 – Title
	%------------------------------------------------
	\begin{frame}
		\centering
		
		{\Large\bfseries Automatic Circuit Discovery for\\[4pt]
			Explainable Network Intrusion Detection}
		
		\vspace{14pt}
		
		{\normalsize
			Adapting ACDC for Mechanistic Interpretability in NIDS}
		
		\vspace{22pt}
		
		{\small
			\textbf{Based on:} Arthur Conmy, Augustine N. Mavor-Parker,\\
			Aengus Lynch, Stefan Heimersheim, Adri\`a Garriga-Alonso}
		
		\vfill
		
		{\small
			Presented by: \textbf{Anbuselvi S}}
	\end{frame}
	
	\begin{frame}{Problem Framing: Interpretable Intrusion Detection}
		\justifying
		Deep learning–based Network Intrusion Detection Systems (NIDS) have become
		central to modern security infrastructures due to their ability to model
		complex temporal and statistical attack patterns. However, these models operate
		as opaque computational systems, providing little insight into \emph{how}
		intrusion decisions are internally computed.
		
		\vspace{6pt}
		From a security operations perspective, the inability to explain alerts
		complicates trust, auditing, and forensic analysis. From a development and
		research perspective, this opacity hinders debugging, robustness analysis,
		and systematic improvement of NIDS models.
		
		\vspace{6pt}
		This work investigates whether \textbf{Automatic Circuit Discovery (ACDC)}
		can be used to expose the internal computational mechanisms responsible for
		intrusion detection decisions.
	\end{frame}
	
	\begin{frame}{Why Existing XAI Is Insufficient for NIDS}
		\begin{itemize}
			\setstretch{1.4}
			\item Popular explainability methods (SHAP, LIME):
			\begin{itemize}
				\item attribute importance to input features
				\item rely on local correlations
			\end{itemize}
			\item These methods do \textbf{not} explain:
			\begin{itemize}
				\item how internal representations interact
				\item which internal components are causally necessary
			\end{itemize}
			\item In security-critical systems:
			\begin{itemize}
				\item correlated features may be spurious
				\item explanations must reflect true decision mechanisms
			\end{itemize}
			\item This motivates a shift from feature attribution to
			\textbf{mechanistic interpretability}.
		\end{itemize}
	\end{frame}
	
	\begin{frame}{Applying ACDC to NIDS: Core Idea}
		\begin{itemize}
			\setstretch{1.4}
			\item ACDC is applied \textbf{post-hoc} to a trained NIDS model
			\item The NIDS model is treated as a computational system, not a black box
			\item The key objective is to identify:
			\begin{itemize}
				\item which internal components
				\item and which information pathways
			\end{itemize}
			are \textbf{causally responsible} for detecting a specific attack
			\item This is achieved by systematically testing and pruning the model’s
			internal computational graph
		\end{itemize}
	\end{frame}
	
	\begin{frame}{Behavior Definition in the Context of NIDS}
		\begin{itemize}
			\setstretch{1.4}
			\item ACDC requires a precise definition of \textbf{behavior}
			\item In NIDS, a behavior corresponds to:
			\begin{itemize}
				\item the detection of a specific attack class
				\item under a defined traffic distribution
			\end{itemize}
			\item Examples of behaviors:
			\begin{itemize}
				\item DDoS traffic classified as malicious
				\item Port scanning patterns flagged as intrusion
			\end{itemize}
			\item Each behavior is quantified using:
			\begin{itemize}
				\item model confidence or logit difference
				\item between clean (attack) and corrupted (benign) traffic
			\end{itemize}
		\end{itemize}
	\end{frame}
	
	\begin{frame}{Model Representation: NIDS as a Computational Graph}
		\begin{itemize}
			\setstretch{1.4}
			\item To apply ACDC, the trained NIDS is converted into a
			\textbf{directed computational graph}
			\item Graph nodes represent:
			\begin{itemize}
				\item feature embeddings
				\item hidden neurons or layers
				\item temporal or attention-based components
			\end{itemize}
			\item Graph edges represent information flow during inference
			\item The intrusion prediction node serves as the behavioral output
			\item ACDC operates by selectively removing edges while preserving behavior
		\end{itemize}
	\end{frame}
	
	\begin{frame}{Clean vs Corrupted Traffic in ACDC for NIDS}
		\begin{itemize}
			\setstretch{1.4}
			\item ACDC relies on comparing clean and corrupted executions
			\item In the context of NIDS:
			\begin{itemize}
				\item \textbf{Clean traffic}: flows containing a specific attack pattern
				\item \textbf{Corrupted traffic}: benign traffic or traffic with attack cues removed
			\end{itemize}
			\item Corruption strategies may include:
			\begin{itemize}
				\item shuffling temporal order
				\item masking protocol-specific features
				\item replacing flow statistics with benign distributions
			\end{itemize}
			\item This setup enables causal testing of internal components
		\end{itemize}
	\end{frame}
	
	\begin{frame}{Activation Patching for Intrusion Detection}
		\begin{itemize}
			\setstretch{1.4}
			\item Activation patching replaces internal activations during inference
			\item For a given component:
			\begin{itemize}
				\item run the model on clean attack traffic
				\item replace selected internal activations with those from corrupted traffic
			\end{itemize}
			\item Measure the change in intrusion prediction
			\item If the prediction remains unchanged:
			\begin{itemize}
				\item the patched component is not causally required
			\end{itemize}
		\end{itemize}
	\end{frame}
	
	\begin{frame}{Recursive Pruning of the Computational Graph}
		\begin{itemize}
			\setstretch{1.4}
			\item ACDC applies activation patching to graph edges
			\item Edges are tested in reverse topological order
			\item For each edge:
			\begin{itemize}
				\item patch the edge using corrupted activations
				\item evaluate the effect on attack detection
			\end{itemize}
			\item If the effect is below a threshold:
			\begin{itemize}
				\item the edge is permanently removed
			\end{itemize}
			\item This process yields a minimal, behavior-preserving circuit
		\end{itemize}
	\end{frame}
	
	\begin{frame}{Interpreting Extracted Circuits in NIDS}
		\begin{itemize}
			\setstretch{1.4}
			\item The extracted circuit represents the minimal internal computation
			\item It captures:
			\begin{itemize}
				\item which features interact
				\item which temporal dependencies matter
				\item which internal representations are essential
			\end{itemize}
			\item Circuits are:
			\begin{itemize}
				\item attack-specific
				\item model-specific
			\end{itemize}
			\item This enables mechanistic understanding of intrusion detection
		\end{itemize}
	\end{frame}
	
	\begin{frame}{What ACDC-Based Interpretability Enables}
		\begin{itemize}
			\setstretch{1.4}
			\item From a security operations perspective:
			\begin{itemize}
				\item improved alert trust
				\item clearer forensic explanations
			\end{itemize}
			\item From a developer perspective:
			\begin{itemize}
				\item identification of spurious correlations
				\item debugging of failure cases
			\end{itemize}
			\item From a research perspective:
			\begin{itemize}
				\item causal analysis of model behavior
				\item principled comparison across attack types
			\end{itemize}
		\end{itemize}
	\end{frame}
	
	\begin{frame}{Example: Circuit Discovery for DDoS Detection}
		\begin{itemize}
			\setstretch{1.4}
			\item Consider a trained NIDS detecting DDoS attacks
			\item Behavior of interest:
			\begin{itemize}
				\item high-confidence prediction for volumetric flooding traffic
			\end{itemize}
			\item ACDC isolates a circuit that:
			\begin{itemize}
				\item aggregates packet rate features
				\item emphasizes short-term temporal bursts
				\item suppresses protocol-specific noise
			\end{itemize}
			\item The extracted circuit explains \emph{how} the model detects DDoS
		\end{itemize}
	\end{frame}
	
	\begin{frame}{Limits of Feature Attribution in NIDS}
		\begin{itemize}
			\setstretch{1.4}
			\item Feature attribution methods explain:
			\begin{itemize}
				\item which inputs influence a decision
			\end{itemize}
			\item They do not explain:
			\begin{itemize}
				\item how features are internally combined
				\item which intermediate representations matter
				\item whether correlations are causal
			\end{itemize}
			\item ACDC instead explains:
			\begin{itemize}
				\item internal computation paths
				\item causal necessity of components
			\end{itemize}
		\end{itemize}
	\end{frame}
	
	\begin{frame}{Evaluating ACDC in Network Intrusion Detection}
		\begin{itemize}
			\setstretch{1.4}
			\item Evaluation focuses on two core questions:
			\begin{itemize}
				\item Does the circuit preserve attack detection performance?
				\item Is the circuit significantly smaller than the full model?
			\end{itemize}
			\item Metrics include:
			\begin{itemize}
				\item prediction divergence (e.g., KL divergence)
				\item detection accuracy under ablation
				\item circuit sparsity
			\end{itemize}
			\item Circuits are validated on held-out attack traffic
		\end{itemize}
	\end{frame}
	
	\begin{frame}{System-Level Impact of ACDC-Enabled NIDS}
		\begin{itemize}
			\setstretch{1.4}
			\item For Security Operations Centers (SOC):
			\begin{itemize}
				\item transparent alerts
				\item reduced false-positive investigation cost
			\end{itemize}
			\item For developers:
			\begin{itemize}
				\item identification of brittle feature dependencies
				\item targeted model refinement
			\end{itemize}
			\item For researchers:
			\begin{itemize}
				\item principled causal analysis of IDS models
				\item comparison of internal mechanisms across attacks
			\end{itemize}
		\end{itemize}
	\end{frame}
	
	\begin{frame}{Research Significance}
		\begin{itemize}
			\setstretch{1.4}
			\item This work extends mechanistic interpretability beyond NLP
			\item Introduces circuit-based explanations for NIDS
			\item Bridges security analytics and causal interpretability
			\item Enables:
			\begin{itemize}
				\item explainable intrusion detection
				\item auditable security decisions
				\item robust and debuggable models
			\end{itemize}
		\end{itemize}
	\end{frame}
	
	\begin{frame}{Limitations of ACDC for NIDS}
		\begin{itemize}
			\setstretch{1.4}
			\item ACDC is computationally intensive
			\item Circuit discovery is performed offline
			\item Results depend on:
			\begin{itemize}
				\item choice of corrupted traffic
				\item behavior definition granularity
			\end{itemize}
			\item Extracted circuits are:
			\begin{itemize}
				\item model-specific
				\item dataset-dependent
			\end{itemize}
		\end{itemize}
	\end{frame}
	
	\begin{frame}{Threat Model and Adversarial Considerations}
		\begin{itemize}
			\setstretch{1.4}
			\item Attackers may attempt to:
			\begin{itemize}
				\item exploit spurious correlations in NIDS
				\item evade detection via adversarial traffic
			\end{itemize}
			\item ACDC enables:
			\begin{itemize}
				\item identification of brittle internal circuits
				\item analysis of attack-specific vulnerabilities
			\end{itemize}
			\item This supports proactive hardening of NIDS models
		\end{itemize}
	\end{frame}
	
	
	\begin{frame}{Deployment Perspective in SOC Environments}
		\begin{itemize}
			\setstretch{1.4}
			\item ACDC is not part of the real-time detection path
			\item It operates as:
			\begin{itemize}
				\item an offline analysis and auditing module
			\end{itemize}
			\item Outputs support:
			\begin{itemize}
				\item analyst trust
				\item incident post-mortem analysis
				\item model debugging and retraining
			\end{itemize}
			\item Compatible with existing SOC pipelines
		\end{itemize}
	\end{frame}
	
	\begin{frame}{Future Research Directions}
		\begin{itemize}
			\setstretch{1.4}
			\item Scaling ACDC to:
			\begin{itemize}
				\item deeper NIDS architectures
				\item multi-attack joint circuits
			\end{itemize}
			\item Integrating gradient-based pruning for efficiency
			\item Studying circuit stability across datasets
			\item Extending mechanistic interpretability to:
			\begin{itemize}
				\item adaptive and evolving attacks
			\end{itemize}
		\end{itemize}
	\end{frame}
	
	
	\begin{frame}{Conclusion}
		\begin{itemize}
			\setstretch{1.4}
			\item Network intrusion detection requires more than accuracy
			\item ACDC enables:
			\begin{itemize}
				\item causal understanding of NIDS decisions
				\item attack-specific internal explanations
			\end{itemize}
			\item This work bridges:
			\begin{itemize}
				\item deep learning security
				\item mechanistic interpretability
			\end{itemize}
			\item Provides a foundation for trustworthy NIDS systems
		\end{itemize}
	\end{frame}	
	
	%------------------------------------------------
	% Thank You Slide
	%------------------------------------------------
	\begin{frame}
		\centering
		\Large{Thank You!}
	\end{frame}
	
\end{document}
